% !TEX program = pdflatex
% =====================================================================
% Zwätschgen-Pubquiz -- Fragenkatalog
% \loesungtrue  -> Version für Quizmaster/Auswertung mit Lösungen
% \loesungfalse -> Fragenheft für die Teams
% =====================================================================
\documentclass[a4paper,11pt]{article}
% =====================================================================
% Gemeinsame Präambel für main.tex (Konzept) und fragenkatalog.tex
% Einbinden mit % =====================================================================
% Gemeinsame Präambel für main.tex (Konzept) und fragenkatalog.tex
% Einbinden mit % =====================================================================
% Gemeinsame Präambel für main.tex (Konzept) und fragenkatalog.tex
% Einbinden mit \input{praeambel}
% =====================================================================
\usepackage[utf8]{inputenc}
\usepackage[T1]{fontenc}
\usepackage{lmodern}
\usepackage[ngerman]{babel}
\usepackage[margin=2.3cm]{geometry}
\usepackage{enumitem}
\usepackage{booktabs}
\usepackage{tabularx}
\usepackage{xcolor}
\usepackage{listings}
\usepackage{fancyhdr}
\usepackage{amsmath}
\usepackage{array}
\usepackage[hidelinks]{hyperref}

% ---------- Farben ----------
\definecolor{zwetschge}{HTML}{5B2A86}
\definecolor{offen}{HTML}{C0392B}

% ---------- Kopf- und Fusszeile ----------
\pagestyle{fancy}
\fancyhf{}
\lhead{Zwätschgenabend -- Pubquiz}
\rhead{HSLU Elektrotechnik}
\cfoot{\thepage}
\setlength{\headheight}{14pt}

% ---------- Listen ----------
\setlist[itemize]{itemsep=0.2em, topsep=0.3em}

% ---------- Hilfsmakros ----------
% Markiert Punkte, die das OK noch klären muss
\newcommand{\offen}[1]{\textcolor{offen}{\textbf{[offen:} #1\textbf{]}}}

\lstset{
  basicstyle=\ttfamily\small,
  frame=single,
  rulecolor=\color{zwetschge!50},
  language=C,
  columns=fullflexible,
  xleftmargin=1em
}

% =====================================================================
\usepackage[utf8]{inputenc}
\usepackage[T1]{fontenc}
\usepackage{lmodern}
\usepackage[ngerman]{babel}
\usepackage[margin=2.3cm]{geometry}
\usepackage{enumitem}
\usepackage{booktabs}
\usepackage{tabularx}
\usepackage{xcolor}
\usepackage{listings}
\usepackage{fancyhdr}
\usepackage{amsmath}
\usepackage{array}
\usepackage[hidelinks]{hyperref}

% ---------- Farben ----------
\definecolor{zwetschge}{HTML}{5B2A86}
\definecolor{offen}{HTML}{C0392B}

% ---------- Kopf- und Fusszeile ----------
\pagestyle{fancy}
\fancyhf{}
\lhead{Zwätschgenabend -- Pubquiz}
\rhead{HSLU Elektrotechnik}
\cfoot{\thepage}
\setlength{\headheight}{14pt}

% ---------- Listen ----------
\setlist[itemize]{itemsep=0.2em, topsep=0.3em}

% ---------- Hilfsmakros ----------
% Markiert Punkte, die das OK noch klären muss
\newcommand{\offen}[1]{\textcolor{offen}{\textbf{[offen:} #1\textbf{]}}}

\lstset{
  basicstyle=\ttfamily\small,
  frame=single,
  rulecolor=\color{zwetschge!50},
  language=C,
  columns=fullflexible,
  xleftmargin=1em
}

% =====================================================================
\usepackage[utf8]{inputenc}
\usepackage[T1]{fontenc}
\usepackage{lmodern}
\usepackage[ngerman]{babel}
\usepackage[margin=2.3cm]{geometry}
\usepackage{enumitem}
\usepackage{booktabs}
\usepackage{tabularx}
\usepackage{xcolor}
\usepackage{listings}
\usepackage{fancyhdr}
\usepackage{amsmath}
\usepackage{array}
\usepackage[hidelinks]{hyperref}

% ---------- Farben ----------
\definecolor{zwetschge}{HTML}{5B2A86}
\definecolor{offen}{HTML}{C0392B}

% ---------- Kopf- und Fusszeile ----------
\pagestyle{fancy}
\fancyhf{}
\lhead{Zwätschgenabend -- Pubquiz}
\rhead{HSLU Elektrotechnik}
\cfoot{\thepage}
\setlength{\headheight}{14pt}

% ---------- Listen ----------
\setlist[itemize]{itemsep=0.2em, topsep=0.3em}

% ---------- Hilfsmakros ----------
% Markiert Punkte, die das OK noch klären muss
\newcommand{\offen}[1]{\textcolor{offen}{\textbf{[offen:} #1\textbf{]}}}

\lstset{
  basicstyle=\ttfamily\small,
  frame=single,
  rulecolor=\color{zwetschge!50},
  language=C,
  columns=fullflexible,
  xleftmargin=1em
}

\usepackage{pgffor}

\newif\ifloesung
\loesungtrue

% ---------- Makros für Fragen ----------
\newcommand{\loesung}[1]{\ifloesung\par\smallskip\noindent\textcolor{zwetschge}{\textbf{Lösung:} #1}\fi}

% Multiple Choice: \mc{Frage}{A}{B}{C}{D}{Lösung mit Erklärung}
\newcommand{\mc}[6]{%
  \item #1\par\nopagebreak\smallskip
  \noindent\begin{tabularx}{\linewidth}{@{}>{\bfseries}l X >{\bfseries}l X@{}}
    A) & #2 & B) & #3 \\[0.25em]
    C) & #4 & D) & #5
  \end{tabularx}%
  \loesung{#6}}

% Schätzfrage: \schaetz{Frage}{Lösung}
\newcommand{\schaetz}[2]{%
  \item #1\par\smallskip\noindent Schätzung: \rule{5cm}{0.4pt}%
  \loesung{#2}}

\newcommand{\runde}[2]{\section{Runde #1: #2}}

\setlist[enumerate]{label=\textbf{\arabic*.}, itemsep=1.1em, topsep=0.6em}

% ---------- Antwortblatt ----------
\newcommand{\kaestchen}{\fbox{\rule{0pt}{0.7em}\rule{0.7em}{0pt}}}
\newcommand{\antwortblatt}[2]{%
  \subsection*{Runde #1}
  \foreach \n in {1,...,#2}{%
    \noindent\makebox[2.5em][l]{\textbf{\n.}}%
    A~\kaestchen\quad B~\kaestchen\quad C~\kaestchen\quad D~\kaestchen\par\medskip}}

\title{\textcolor{zwetschge}{\textbf{Zwätschgen-Pubquiz}}\\[0.3em]\large Fragenkatalog\ifloesung{} mit Lösungen\fi}
\author{OK Zwätschgenabend, HSLU Elektrotechnik}
\date{\today}

\begin{document}
\maketitle
\thispagestyle{fancy}

\noindent Jede Multiple-Choice-Frage hat genau eine richtige Antwort. Die Fragen sind bewusst schwer, aber mit Überlegen, Rechnen oder Ausschlussverfahren lösbar. Rechnen im Kopf oder auf dem Antwortblatt ist erlaubt, Handys nicht.

% =====================================================================
\runde{1}{Elektrotechnik mit Tücken}
\begin{enumerate}
  \mc{Wie gross ist der Effektivwert einer symmetrischen Rechteckspannung, die zwischen $+10\,\mathrm{V}$ und $-10\,\mathrm{V}$ wechselt?}
     {$7{,}07\,\mathrm{V}$}{$5{,}00\,\mathrm{V}$}{$10\,\mathrm{V}$}{$6{,}37\,\mathrm{V}$}
     {C. Der Betrag ist immer $10\,\mathrm{V}$, also ist auch der Effektivwert $10\,\mathrm{V}$. $7{,}07\,\mathrm{V}$ gilt nur für den Sinus.}

  \mc{Und bei einer Dreieckspannung mit Scheitelwert $10\,\mathrm{V}$?}
     {$7{,}07\,\mathrm{V}$}{$5{,}77\,\mathrm{V}$}{$5{,}00\,\mathrm{V}$}{$6{,}37\,\mathrm{V}$}
     {B. $\hat{u}/\sqrt{3} \approx 5{,}77\,\mathrm{V}$.}

  \mc{Ein Kondensator $C$ ist auf $U$ geladen und wird mit einem gleich grossen, ungeladenen Kondensator verbunden. Wie viel der ursprünglichen Energie steckt danach noch in den beiden Kondensatoren?}
     {$100\,\%$}{$75\,\%$}{$25\,\%$}{$50\,\%$}
     {D. Die Ladung verteilt sich, beide haben $U/2$. Energie: $2 \cdot \tfrac{1}{2}C(U/2)^2 = \tfrac{1}{4}CU^2$, also die Hälfte. Der Rest geht im Widerstand oder Funken verloren, unabhängig davon, wie klein er ist.}

  \mc{Durch eine Spule mit $10\,\mathrm{mH}$ fliessen $2\,\mathrm{A}$. Der Strom wird in $1\,\mathrm{\mu s}$ abgeschaltet. Welche Spannung entsteht ungefähr?}
     {$20\,\mathrm{kV}$}{$2\,\mathrm{kV}$}{$200\,\mathrm{V}$}{$20\,\mathrm{V}$}
     {A. $u = L\,\Delta i/\Delta t = 10^{-2} \cdot 2 / 10^{-6} = 20\,\mathrm{kV}$. Deshalb gehört eine Freilaufdiode über jede Relaisspule.}

  \mc{Ein Kupferdraht wird durch einen ersetzt, der doppelt so lang ist und den doppelten Durchmesser hat. Was passiert mit dem Widerstand?}
     {Er verdoppelt sich}{Er bleibt gleich}{Er halbiert sich}{Er viertelt sich}
     {C. $R \sim l/A$: doppelte Länge und vierfache Fläche ergibt $2/4 = 1/2$.}

  \mc{Wie gross ist der Wirkungsgrad bei Leistungsanpassung ($R_L = R_i$)?}
     {$100\,\%$}{$50\,\%$}{$71\,\%$}{$25\,\%$}
     {B. Die Hälfte der Leistung wird im Innenwiderstand verheizt. Maximale Leistung heisst nicht maximaler Wirkungsgrad.}

  \mc{Ein Kupferleiter erwärmt sich um $100\,\mathrm{K}$. Um wie viel steigt sein Widerstand ungefähr?}
     {$0{,}4\,\%$}{$4\,\%$}{$40\,\%$}{$400\,\%$}
     {C. $\alpha_{Cu} \approx 0{,}0039\,\mathrm{K^{-1}}$, also etwa $+39\,\%$.}

  \mc{Ein Widerstand hat die Farbringe Braun -- Schwarz -- Orange -- Gold. Welchen Wert hat er?}
     {$1\,\mathrm{k\Omega} \pm 5\,\%$}{$10\,\mathrm{k\Omega} \pm 10\,\%$}{$100\,\mathrm{k\Omega} \pm 5\,\%$}{$10\,\mathrm{k\Omega} \pm 5\,\%$}
     {D. 1, 0, $\times 10^3$, Gold $= \pm 5\,\%$ (Silber wäre $\pm 10\,\%$).}

  \mc{Eine $100$-W- und eine $25$-W-Glühlampe (beide für $230\,\mathrm{V}$) werden in Serie an $230\,\mathrm{V}$ angeschlossen. Welche leuchtet heller?}
     {Die 25-W-Lampe}{Die 100-W-Lampe}{Beide gleich hell}{Keine leuchtet sichtbar}
     {A. Die 25-W-Lampe hat den viermal grösseren Widerstand. Bei gleichem Strom fällt an ihr die grössere Leistung $P = I^2R$ ab.}
\end{enumerate}

% =====================================================================
\runde{2}{Digitaltechnik und Embedded}
\begin{enumerate}
  \mc{\texttt{int8\_t x = 127; x++;} Welchen Wert hat \texttt{x} danach auf einem üblichen Compiler?}
     {128}{0}{$-128$}{$-127$}
     {C. Überlauf im Zweierkomplement: \texttt{0x7F} $+ 1 =$ \texttt{0x80} $= -128$.}

  \mc{Wie viele verschiedene Boolesche Funktionen mit zwei Eingängen gibt es?}
     {4}{8}{256}{16}
     {D. Eine Wahrheitstabelle mit 4 Zeilen, jede Zeile 0 oder 1: $2^4 = 16$.}

  \mc{Wie lautet $-1$ als 8-Bit-Zahl im Zweierkomplement?}
     {\texttt{0xFF}}{\texttt{0x81}}{\texttt{0x80}}{\texttt{0xFE}}
     {A. Alle Bits gesetzt. \texttt{0x81} wäre $-1$ in Vorzeichen-Betrag-Darstellung.}

  \mc{Mit welchem Gatter allein lässt sich jede beliebige Logikfunktion aufbauen?}
     {XOR}{NOR}{AND}{OR}
     {B. NOR (wie NAND) ist universell.}

  \mc{Warum brauchen die I\textsuperscript{2}C-Leitungen SDA und SCL Pull-up-Widerstände?}
     {Zur Impedanzanpassung}{Als ESD-Schutz}{Weil die Ausgänge Open-Drain sind}{Zur Taktsynchronisation}
     {C. Die Teilnehmer ziehen die Leitung nur auf Low. High entsteht über den Pull-up. Das erlaubt auch Clock-Stretching und Arbitrierung.}

  \mc{Eine UART läuft mit $115\,200\,\mathrm{Baud}$ im Format 8N1. Wie viele Nutzbytes pro Sekunde sind maximal möglich?}
     {14\,400}{115\,200}{12\,800}{11\,520}
     {D. Pro Byte werden 10 Bit übertragen (Start, 8 Daten, Stopp).}

  \mc{Ein 32-Bit-Zähler wird mit $1\,\mathrm{MHz}$ getaktet. Nach welcher Zeit läuft er über?}
     {$4{,}3\,\mathrm{s}$}{$71{,}6\,\mathrm{min}$}{$49{,}7\,\mathrm{Tage}$}{$136\,\mathrm{Jahre}$}
     {B. $2^{32}\,\mathrm{\mu s} \approx 4295\,\mathrm{s}$. 49,7 Tage wären es bei einem Millisekunden-Tick (der bekannte \texttt{millis()}-Überlauf).}

  \mc{Was bewirkt das Schlüsselwort \texttt{volatile} in C?}
     {Der Compiler darf Zugriffe nicht wegoptimieren}{Die Variable liegt im RAM statt im Flash}{Die Variable ist atomar und threadsicher}{Die Variable wird beim Reset nicht initialisiert}
     {A. Jeder Zugriff wird wirklich ausgeführt (wichtig bei Registern und ISR-Variablen). Atomar wird sie dadurch nicht.}

  \mc{Ein asynchrones Signal wird über zwei hintereinander geschaltete Flip-Flops in eine Taktdomäne eingelesen. Wozu?}
     {Zum Entprellen eines Tasters}{Zur Frequenzhalbierung}{Um Metastabilität abzufangen}{Zur Unterdrückung von Glitches}
     {C. Das erste Flip-Flop darf metastabil werden. Bis zur nächsten Taktflanke hat es sich mit hoher Wahrscheinlichkeit eingeschwungen.}

  \mc{Wie viele Adressleitungen braucht ein byteadressierter Speicher mit $64\,\mathrm{KiB}$?}
     {20}{32}{64}{16}
     {D. $64 \cdot 1024 = 2^{16}$.}
\end{enumerate}

% =====================================================================
\runde{3}{Signale, Regelung, Kommunikation}
\begin{enumerate}
  \mc{Welche Harmonischen enthält eine symmetrische Rechteckschwingung (ohne Gleichanteil)?}
     {Nur gerade}{Nur ungerade}{Alle}{Nur die Grundschwingung}
     {B. $f_0, 3f_0, 5f_0, \dots$ mit Amplituden $\sim 1/n$.}

  \mc{Die Spannung an einem Widerstand wird verdoppelt. Um wie viel steigt der Pegel?}
     {$+3\,\mathrm{dB}$}{$+10\,\mathrm{dB}$}{$+6\,\mathrm{dB}$}{$+20\,\mathrm{dB}$}
     {C. $20\log_{10}2 \approx 6\,\mathrm{dB}$. $3\,\mathrm{dB}$ gelten für die doppelte Leistung.}

  \mc{Welcher Leistung entspricht $0\,\mathrm{dBm}$?}
     {$1\,\mathrm{mW}$}{$1\,\mathrm{W}$}{$1\,\mathrm{\mu W}$}{$0\,\mathrm{W}$}
     {A. dBm ist auf $1\,\mathrm{mW}$ bezogen.}

  \mc{Worauf reagiert der D-Anteil eines PID-Reglers?}
     {Auf den Betrag der Regelabweichung}{Auf das Integral der Regelabweichung}{Auf den Sollwert}{Auf die Änderungsgeschwindigkeit der Regelabweichung}
     {D. Deshalb verstärkt er auch Messrauschen stark.}

  \mc{Ein System hat einen Pol in der rechten $s$-Halbebene. Was bedeutet das?}
     {Es schwingt gedämpft}{Es ist grenzstabil}{Es ist instabil}{Es reagiert schneller}
     {C. Der Pol führt zu einem aufklingenden Exponentialterm $e^{\sigma t}$ mit $\sigma > 0$.}

  \mc{Welche Wellenlänge hat ein WLAN-Signal bei $2{,}4\,\mathrm{GHz}$?}
     {$1{,}25\,\mathrm{cm}$}{$12{,}5\,\mathrm{cm}$}{$1{,}25\,\mathrm{m}$}{$12{,}5\,\mathrm{m}$}
     {B. $\lambda = c/f = 3\cdot10^8/2{,}4\cdot10^9 = 0{,}125\,\mathrm{m}$.}

  \mc{Wie viele Bit überträgt ein Symbol bei 64-QAM?}
     {6}{4}{8}{64}
     {A. $\log_2 64 = 6$.}

  \mc{Einer Faltung im Zeitbereich entspricht im Frequenzbereich eine \dots}
     {Addition}{Faltung}{Ableitung}{Multiplikation}
     {D. Faltungssatz der Fourier-Transformation.}

  \mc{Mit welcher Steigung fällt der Amplitudengang eines idealen Integrators im Bode-Diagramm?}
     {$-40\,\mathrm{dB/Dekade}$}{$-20\,\mathrm{dB/Dekade}$}{$-6\,\mathrm{dB/Dekade}$}{$-3\,\mathrm{dB/Dekade}$}
     {B. Das entspricht $-6\,\mathrm{dB}$ pro \emph{Oktave}, daher die Falle bei C.}

  \mc{Wie sieht das Spektrum eines idealen Dirac-Impulses aus?}
     {Eine Sinc-Funktion}{Ein einzelner Strich bei $0\,\mathrm{Hz}$}{Konstant über alle Frequenzen}{Eine Gausskurve}
     {C. Der Dirac enthält alle Frequenzen gleich stark. Der Strich bei $0\,\mathrm{Hz}$ gehört zum Gleichsignal.}
\end{enumerate}

% =====================================================================
\runde{4}{Energie und Schweiz}
\noindent\textit{Hinweis für das OK: Zeitabhängige Werte vor dem Anlass prüfen.}
\begin{enumerate}
  \mc{Mit welcher Frequenz arbeitet das Bahnstromnetz der SBB?}
     {$50\,\mathrm{Hz}$}{$25\,\mathrm{Hz}$}{$16{,}7\,\mathrm{Hz}$}{$60\,\mathrm{Hz}$}
     {C. $16{,}7\,\mathrm{Hz}$ ($16\,\tfrac{2}{3}\,\mathrm{Hz}$) bei $15\,\mathrm{kV}$ Fahrleitungsspannung.}

  \mc{Welches ist die höchste Spannungsebene im Schweizer Übertragungsnetz?}
     {$220\,\mathrm{kV}$}{$380\,\mathrm{kV}$}{$500\,\mathrm{kV}$}{$750\,\mathrm{kV}$}
     {B. $380\,\mathrm{kV}$ (dazu die 220-kV-Ebene).}

  \mc{Welche Quelle liefert den grössten Anteil der Schweizer Stromproduktion?}
     {Kernenergie}{Photovoltaik}{Windkraft}{Wasserkraft}
     {D. Wasserkraft, deutlich über die Hälfte.}

  \mc{Wer betreibt das Schweizer Übertragungsnetz?}
     {Swissgrid}{ElCom}{Axpo}{BKW}
     {A. Swissgrid. Die ElCom ist die Regulierungsbehörde.}

  \mc{Bei welcher Netzfrequenz beginnt im europäischen Verbundnetz die erste Stufe des automatischen Lastabwurfs?}
     {$49{,}8\,\mathrm{Hz}$}{$49{,}5\,\mathrm{Hz}$}{$47{,}5\,\mathrm{Hz}$}{$49{,}0\,\mathrm{Hz}$}
     {D. Bei $49{,}0\,\mathrm{Hz}$. Bei $47{,}5\,\mathrm{Hz}$ trennen sich die Kraftwerke vom Netz.}

  \mc{Welche Leistung hat das Pumpspeicherkraftwerk Nant de Drance im Wallis?}
     {$90\,\mathrm{MW}$}{$900\,\mathrm{MW}$}{$9\,\mathrm{GW}$}{$300\,\mathrm{MW}$}
     {B. $900\,\mathrm{MW}$, etwa so viel wie ein grosses Kernkraftwerk.}

  \mc{Welchen Wirkungsgrad haben heutige monokristalline PV-Module typischerweise?}
     {$8$--$10\,\%$}{$35$--$40\,\%$}{$20$--$23\,\%$}{$50$--$55\,\%$}
     {C. Etwa $20$--$23\,\%$. Werte um $40\,\%$ erreichen nur Mehrfachzellen im Labor oder in der Raumfahrt.}

  \mc{An wie vielen Standorten sind in der Schweiz Kernkraftwerke in Betrieb?}
     {2}{4}{5}{3}
     {D. Drei Standorte: Beznau, Gösgen und Leibstadt. Zusammen sind es 4 Reaktoren, daher die Falle bei B. Mühleberg ist seit 2019 abgeschaltet.}

  \mc{Wie heisst der Schweizer Steckertyp nach internationaler Bezeichnung?}
     {Typ J}{Typ F}{Typ G}{Typ C}
     {A. Typ J (T12/T13). Typ F ist der Schuko-Stecker, Typ G der britische.}
\end{enumerate}

% =====================================================================
\runde{5}{Geschichte und Köpfe}
\begin{enumerate}
  \mc{Wer kämpfte im «Stromkrieg» Ende des 19.\ Jahrhunderts für den Gleichstrom?}
     {Nikola Tesla}{George Westinghouse}{Thomas Edison}{Werner von Siemens}
     {C. Edison. Tesla und Westinghouse setzten sich mit Wechselstrom durch.}

  \mc{An welcher Hochschule war Konrad Zuses Rechner Z4 ab 1950 in Betrieb?}
     {TU München}{ETH Zürich}{EPFL Lausanne}{University of Cambridge}
     {B. An der ETH Zürich, damals einer der ersten kommerziell genutzten Computer Europas.}

  \mc{Welche Programmiersprache entwickelte Niklaus Wirth an der ETH Zürich?}
     {C}{Java}{Python}{Pascal}
     {D. Pascal (später auch Modula-2 und Oberon).}

  \mc{Wo wurde das World Wide Web erfunden?}
     {Am CERN}{Am MIT}{In den Bell Labs}{Bei Xerox PARC}
     {A. Am CERN in Genf, von Tim Berners-Lee (1989).}

  \mc{Wer erhielt als einzige Person zweimal den Nobelpreis für Physik?}
     {William Shockley}{Albert Einstein}{Richard Feynman}{John Bardeen}
     {D. John Bardeen: 1956 für den Transistor, 1972 für die Theorie der Supraleitung.}

  \mc{In welchem heutigen Land wurde Nikola Tesla geboren?}
     {Serbien}{Kroatien}{Österreich}{Ungarn}
     {B. In Smiljan, im heutigen Kroatien.}

  \mc{Wofür steht das «S» in LASER?}
     {Stimulated}{Spontaneous}{Single}{Spectral}
     {A. Light Amplification by \emph{Stimulated} Emission of Radiation.}

  \mc{Was fand das Team um Grace Hopper 1947 im Relaisrechner Harvard Mark II?}
     {Einen Kaffeefleck}{Eine Kakerlake}{Eine Motte}{Eine Maus}
     {C. Eine Motte. Sie wurde ins Logbuch geklebt, als «first actual case of bug being found».}

  \mc{Welche Datenbreite hatte der erste Mikroprozessor, der Intel 4004 (1971)?}
     {8 Bit}{4 Bit}{16 Bit}{2 Bit}
     {B. 4 Bit, ursprünglich für einen Tischrechner entwickelt.}
\end{enumerate}

% =====================================================================
\runde{6}{Schätzfragen}
\noindent Das nächstliegende Team gewinnt. Bei grossen Zahlen zählt der Faktor, nicht die Differenz.
\begin{enumerate}
  \schaetz{Wie viele Elektronen fliessen bei $1\,\mathrm{A}$ pro Sekunde durch einen Leiterquerschnitt?}
     {ca.\ $6{,}24 \cdot 10^{18}$ $\left(1/1{,}602 \cdot 10^{-19}\right)$}

  \schaetz{Wie schnell bewegen sich die Elektronen (Driftgeschwindigkeit) in einem Kupferdraht mit $1\,\mathrm{mm^2}$ Querschnitt bei $1\,\mathrm{A}$? Angabe in mm/s.}
     {ca.\ $0{,}07\,\mathrm{mm/s}$, also etwa $26\,\mathrm{cm}$ pro Stunde. $v = I/(nAe)$ mit $n \approx 8{,}5 \cdot 10^{28}\,\mathrm{m^{-3}}$.}

  \schaetz{Wie viele Liter Wasser muss man $1\,\mathrm{m}$ hoch pumpen, um $1\,\mathrm{Wh}$ zu speichern?}
     {ca.\ $367\,\mathrm{l}$ ($3600\,\mathrm{J} / (9{,}81\,\mathrm{m/s^2} \cdot 1\,\mathrm{m})$)}

  \schaetz{Wie viele Transistoren hatte der Intel 4004?}
     {ca.\ 2300}

  \schaetz{Wie lang ist der Gotthard-Basistunnel?}
     {ca.\ $57\,\mathrm{km}$}

  \schaetz{Wie hoch ist die Staumauer der Grande Dixence?}
     {$285\,\mathrm{m}$}

  \schaetz{Wie viele Sekunden braucht das Licht von der Sonne zur Erde?}
     {ca.\ $500\,\mathrm{s}$ (8 Minuten 20 Sekunden)}

  \schaetz{Wie viele Kilometer Leitungen umfasst das Schweizer Übertragungsnetz von Swissgrid?}
     {ca.\ $6700\,\mathrm{km}$ (Stand vor dem Anlass prüfen)}

  \schaetz{Um welchen Faktor ist die Energiedichte (pro kg) von Benzin grösser als die eines Li-Ionen-Akkus?}
     {ca.\ 50 (Benzin $\approx 12\,\mathrm{kWh/kg}$, Akkuzelle $\approx 0{,}25\,\mathrm{kWh/kg}$). Antworten zwischen 40 und 60 gelten als sehr gut.}
\end{enumerate}

% =====================================================================
\runde{7}{Zwätschgen-Spezial}
\begin{enumerate}
  \mc{Zu welcher Pflanzenfamilie gehört die Zwetschge?}
     {Nachtschattengewächse}{Rosengewächse}{Lorbeergewächse}{Kürbisgewächse}
     {B. Rosengewächse, wie Apfel, Kirsche und Aprikose.}

  \mc{Was stimmt über Zwetschge und Pflaume?}
     {Die Zwetschge ist eine Unterart der Pflaume}{Sie gehören zu verschiedenen Gattungen}{Die Zwetschge ist eine Kreuzung aus Pflaume und Aprikose}{Es ist dieselbe Frucht, nur der Name ist regional}
     {A. \emph{Prunus domestica} subsp.\ \emph{domestica}. Zwetschgen sind länglicher und haben einen leicht lösbaren Stein.}

  \mc{Was ist der weissliche Belag auf frischen Zwetschgen?}
     {Ein Pestizidrest}{Hefe}{Zuckerkristalle}{Eine natürliche Wachsschicht}
     {D. Eine natürliche Wachsschicht, die die Frucht vor dem Austrocknen schützt.}

  \mc{Wie heisst die getrocknete Zwetschge auf Englisch?}
     {plum}{sloe}{prune}{damson}
     {C. Prune. Sloe ist die Schlehe, damson eine Pflaumensorte.}

  \mc{Welcher Stoff im Zwetschgenkern kann im Körper Blausäure freisetzen?}
     {Solanin}{Amygdalin}{Koffein}{Oxalsäure}
     {B. Amygdalin. Solanin steckt in grünen Kartoffeln.}

  \mc{Man steckt eine Kupfer- und eine Zinkelektrode in eine Zwetschge. Welche Leerlaufspannung misst man ungefähr?}
     {$0{,}1\,\mathrm{V}$}{$1{,}5\,\mathrm{V}$}{$3{,}7\,\mathrm{V}$}{$0{,}9\,\mathrm{V}$}
     {D. Etwa $0{,}9$ bis $1\,\mathrm{V}$. Die Spannung kommt vom Elektrodenpaar, nicht von der Frucht. Live-Demo mit Multimeter möglich.}
\end{enumerate}

% =====================================================================
\section{Finale}
\subsection*{Finalfrage mit Punkteinsatz}
\begin{enumerate}
  \item Bringe diese Erfindungen in die richtige zeitliche Reihenfolge: Mikroprozessor (Intel 4004), Voltasche Säule, integrierter Schaltkreis, Transistor, Glühlampe (Edison).
  \loesung{Voltasche Säule (1800), Glühlampe (1879), Transistor (1947), integrierter Schaltkreis (1958), Mikroprozessor (1971).}
\end{enumerate}

\subsection*{Stechfrage bei Gleichstand}
\begin{enumerate}
  \schaetz{Wie viele Meter Kupferdraht mit $1\,\mathrm{mm^2}$ Querschnitt ergeben einen Widerstand von $1\,\Omega$?}
     {ca.\ $58\,\mathrm{m}$ ($l = RA/\rho$ mit $\rho_{Cu} \approx 0{,}0172\,\Omega\,\mathrm{mm^2/m}$)}
\end{enumerate}

% =====================================================================
\newpage
\section*{Antwortblatt}
\noindent Teamname: \rule{9cm}{0.4pt}

\medskip
\noindent\begin{minipage}[t]{0.48\textwidth}
\antwortblatt{1}{9}
\antwortblatt{2}{10}
\end{minipage}\hfill
\begin{minipage}[t]{0.48\textwidth}
\antwortblatt{3}{10}
\antwortblatt{4}{9}
\end{minipage}

\newpage
\noindent Teamname: \rule{9cm}{0.4pt}

\medskip
\noindent\begin{minipage}[t]{0.48\textwidth}
\antwortblatt{5}{9}
\antwortblatt{7}{6}
\end{minipage}\hfill
\begin{minipage}[t]{0.48\textwidth}
\subsection*{Runde 6 (Schätzfragen)}
\foreach \n in {1,...,9}{\noindent\makebox[2.5em][l]{\textbf{\n.}}\rule{5cm}{0.4pt}\par\bigskip}
\subsection*{Finale}
\noindent Einsatz: \rule{3cm}{0.4pt}\par\bigskip
\foreach \n in {1,...,5}{\noindent\makebox[2.5em][l]{\textbf{\n.}}\rule{5cm}{0.4pt}\par\bigskip}
\end{minipage}

\end{document}
